\graphicspath{{abstracts/workshops/03-sparql-queries/img/}{img/}}

\begin{beitrag}[Rita Reuter, Stefan Schubert]{SPARQL-Queries für Geisteswissenschaftler*innen}{%
Rita Reuter\,\orcidlink{0000-0003-8901-2345} \\
\small \href{mailto:rita.reuter@example.org}{rita.reuter@example.org} \\
\small Deutsches Literaturarchiv Marbach \\[0.5cm]
Stefan Schubert\,\orcidlink{0000-0001-9012-3456} \\
\small \href{mailto:stefan.schubert@example.org}{stefan.schubert@example.org} \\
\small Johannes Gutenberg-Universität Mainz
}

\begin{abstract}
Der Workshop vermittelt Geisteswissenschaftler*innen ohne Vorerfahrung die handwerklichen Grundlagen von SPARQL für WissKI-basierte Forschungsdaten: In vier problemorientierten Modulen werden ausgehend von konkreten Forschungsfragen SELECT-Abfragen, Filterung und Aggregation (einschließlich geisteswissenschaftlicher Datierungen) sowie die Verknüpfung mit externen Quellen über SERVICE und OPTIONAL erarbeitet. Als Arbeitsumgebung dient die frei verfügbare YASGUI-Oberfläche an öffentlich zugänglichen Endpunkten; Ziel ist das konzeptuelle Verständnis dessen, was eine SPARQL-Abfrage leistet und worin sie sich von einer Volltextsuche unterscheidet -- nicht die Beherrschung jedes einzelnen Sprachkonstrukts.
\end{abstract}

\section{Worum es geht}
SPARQL, die Abfragesprache des semantischen Webs, ist die natürliche Schnittstelle zu jeder WissKI-Datenbasis. Sie erlaubt, die in einem Triplestore abgelegten Forschungsdaten frei und ausdrucksstark zu befragen -- vorausgesetzt, man kennt die wichtigsten Sprachkonstrukte und versteht, wie ein RDF-Graph zu \glqq lesen\grqq{} ist. Der vorliegende Workshop wendet sich an Geisteswissenschaftler*innen ohne nennenswerte Vorerfahrung in Abfragesprachen und vermittelt in vier Modulen die handwerklichen Grundlagen, die für die Arbeit mit einem WissKI-basierten SPARQL-Endpunkt erforderlich sind.

\section{Didaktischer Ansatz}
Der Workshop folgt einem ausgesprochen problemorientierten Ansatz: Wir starten nicht mit der Grammatik von SPARQL, sondern mit einer konkreten Forschungsfrage. \glqq Welche Personen sind nach 1750 geboren und stehen mit dem Brief X in Beziehung?\grqq{} ist eine typische Einstiegsfrage, die wir im Verlauf der ersten Stunde gemeinsam in eine SPARQL-Abfrage übersetzen. Aus diesem konkreten Ausgangspunkt heraus arbeiten wir uns systematisch zu den allgemeinen Sprachkonstrukten vor. Die didaktische Begründung dieses Ansatzes ist gesondert dargelegt.\footcite[46]{reuter2024-sparql-didaktik}

\section{Inhalte}
Das erste Modul vermittelt die grundlegenden Strukturen einer SPARQL-Abfrage: \texttt{SELECT}, \texttt{WHERE}, Tripel-Muster, Variablen, Präfix-Deklarationen. Wir arbeiten dabei direkt am Endpunkt eines real existierenden, öffentlich zugänglichen WissKI-Projekts und sehen unsere Abfrage-Ergebnisse in Sekundenschnelle in einer tabellarischen Auswertung.

Das zweite Modul behandelt Filtern, Sortieren und Aggregation: \texttt{FILTER}, \texttt{ORDER BY}, \texttt{LIMIT}, \texttt{GROUP BY} und \texttt{COUNT}. Wir vertiefen die Modul-Inhalte an einer Reihe von Forschungsfragen, deren Beantwortung jeweils ein neues Sprachkonstrukt erforderlich macht. Besonders ausführlich behandeln wir den Umgang mit Datierungen, weil dieser Aspekt in geisteswissenschaftlichen Datenbeständen erfahrungsgemäß sehr schwierig ist.

Das dritte Modul widmet sich der Verbindung mehrerer Datenquellen über die SPARQL-Konstrukte \texttt{SERVICE} und \texttt{OPTIONAL}. Wir reichern unsere lokalen WissKI-Daten mit Informationen aus Wikidata und der GND an und sehen unmittelbar, welche Mehrwerte sich aus der Vernetzung ergeben -- und welche Fallstricke zu beachten sind.

Das vierte Modul ist der eigenen Forschungsfrage gewidmet: Die Teilnehmer*innen bringen eine eigene Frage ein und entwickeln gemeinsam mit den Trainer*innen eine konkrete SPARQL-Lösung. Das Format soll die Teilnehmer*innen in die Lage versetzen, nach dem Workshop selbständig weiterzuarbeiten. Ein Lehrbuch, das den im Workshop vermittelten Stoff vollständig abdeckt, ist als begleitendes Lesematerial zu empfehlen.\footcite[112]{schubert2025-sparql-lehrbuch}

\section{Werkzeugumgebung}
Wir arbeiten mit der frei verfügbaren YASGUI-Oberfläche, die die wichtigsten Komfortfunktionen (Syntax-Hervorhebung, Autovervollständigung, Ergebnisanzeige in mehreren Formaten) bietet, ohne dass eine lokale Installation erforderlich ist. Alle im Workshop verwendeten Endpunkte sind öffentlich zugänglich, sodass die Teilnehmer*innen die Inhalte des Workshops nach Belieben fortsetzen können.

\section{Ein erstes Beispiel}
Zur Veranschaulichung sei hier die Eingangsabfrage des ersten Moduls aufgeführt -- eine einfache Anfrage, die alle Personen aus einer Demosammlung zusammen mit ihrem Sterbedatum zurückgibt.

\begin{lstlisting}[language=SQL, caption={Eine erste SPARQL-Abfrage als Einstieg.}]
PREFIX crm: <http://www.cidoc-crm.org/cidoc-crm/>
PREFIX rdfs: <http://www.w3.org/2000/01/rdf-schema#>

SELECT ?person ?label ?date WHERE {
  ?person a crm:E21_Person ;
          rdfs:label ?label ;
          crm:P100i_died_in/crm:P4_has_time-span/rdfs:label ?date .
}
ORDER BY ?date
\end{lstlisting}

\section{Ergebnis}
Am Ende des Workshops sind die Teilnehmer*innen in der Lage, einfache und mittelkomplexe SPARQL-Abfragen selbstständig zu formulieren, sie schrittweise zu verbessern und ihre Ergebnisse kritisch zu interpretieren. Wichtiger als die Beherrschung jedes einzelnen Sprachkonstrukts ist uns dabei das Verständnis dafür, was eine SPARQL-Abfrage \emph{tut} und worin sie sich von einer Volltextsuche unterscheidet.

\end{beitrag}

\section*{Kurzbiografie(n)}

\subsection*{Rita Reuter}
Rita Reuter ist Mitarbeiterin der Abteilung Digital Humanities am Deutschen Literaturarchiv Marbach. Sie betreut SPARQL-basierte Recherche-Werkzeuge der Marbacher Forschungsdaten und entwickelt didaktische Konzepte zur Vermittlung von Abfragesprachen für Geisteswissenschaftler*innen. Sie ist Autorin mehrerer einschlägiger Aufsätze zur Didaktik der digitalen Geisteswissenschaften.

\subsection*{Stefan Schubert}
Stefan Schubert ist Wissenschaftler an der Johannes Gutenberg-Universität Mainz und Autor eines im erscheinenden Lehrbuchs zur Abfragesprache SPARQL für die Geisteswissenschaften. Er hält regelmäßig Lehrveranstaltungen zu semantischen Datenmodellen und engagiert sich in der Entwicklung offener Lehrmaterialien für die digitalen Geisteswissenschaften.
